% Archival extracts moved out of the scientific supplement. Not a standalone manuscript.
\section{\edit{Candidate settings and numerical protection}}\label{sec:implementation}

\subsection{Raw forecast rules}
\edit{In Table~\ref{tab:expertconfig}, $f_{t,i}$ denotes the supplied forecast.
Windows and lags count observation steps; settings are fixed across datasets and runs.
SD denotes standard deviation.}
We estimate $v_i$ from observations strictly before calibration, using
the population-variance definition. To keep it positive, we apply a floor
equal to the larger of $10^{-12}$ and $10^{-6}$ times the median variance
above $10^{-12}$. If no variance exceeds $10^{-12}$, the floor is
$10^{-12}$.
Lagged observations may precede calibration. Raw-error and feature
histories start at calibration and continue through scoring.

\begin{table}[htbp]\centering\color{red}\small
\caption{Raw forecast candidate settings.}\label{tab:expertconfig}
\begin{tabular}{@{}p{.20\linewidth}p{.75\linewidth}@{}}\toprule
Rule (columns) & Definition and fixed settings\\\midrule
Raw (1) & The supplied forecast $f_{t,i}$.\\
Bias correction (4) & Four candidates, with windows of 3, 7, 14 and 28 observation steps. Each subtracts the mean returned base-forecast error $f_{s,i}-y_{s,i}$ over its window. Use all available errors when fewer have returned; use zero correction initially.\\
Persistence (1) & The most recent observation, $y_{t-1,i}$.\\
Recent mean (1) & Mean of the preceding 7 observations, smoothing recent level fluctuations.\\
Lagged value (1) & The observation 7 steps earlier, $y_{t-7,i}$.\\
Ridge correction (1) & Estimate normalized base-forecast error $(f_{s,i}-y_{s,i})/\sqrt{v_i}$ from up to 36 returned observations; start fitting after 12 errors have returned. Features: current base forecast, most recent observation, preceding 7-observation mean, position within the 12-step prediction block divided by 12, and sine/cosine of the timestamp's calendar weekday. Standardize each feature within the node's fitting window using its mean and population SD (scale one for a constant feature); apply those same statistics to the current feature vector. Use ridge penalty 10 and an unpenalized intercept. Subtract the estimated error multiplied by $\sqrt{v_i}$; before fitting starts, use the raw forecast.\\\bottomrule
\end{tabular}\end{table}

\subsection{Projection and release acceptance}\label{sec:acceptance}
\edit{To reduce rounding errors at the boundary, projection uses intervals
slightly contracted toward the reference:}
\begin{equation}
L^{\rm num}_{t,i}=b_{t,i}+(1-10^{-6})(L_{t,i}-b_{t,i}),\qquad
U^{\rm num}_{t,i}=b_{t,i}+(1-10^{-6})(U_{t,i}-b_{t,i}).
\label{eq:intervalinset}
\end{equation}
\edit{The $10^{-6}$ margin is a fixed numerical setting. Every released
forecast is checked against the original protection envelope. Where coordinate
repair is used, we clip to the contracted intervals and move any entry
that fails the exact envelope test toward the reference, one representable
floating-point step at a time, for at most 16 steps per entry.
We retain the underlying representation only if the repaired forecast
passes both coherence and envelope checks. Otherwise, the release guard
contracts the forecast and its representation jointly toward the reference,
with the reference as the final fallback.}

Clarabel solves the projection quadratic program. Sequential
least-squares programming (SLSQP) supplies fallback solutions when needed.

\begin{table}[H]\centering\small
\caption{Projection and publication settings.}\label{tab:numerics}
\begin{tabular}{@{}p{.23\linewidth}p{.72\linewidth}@{}}\toprule
Operation & Acceptance criteria\\\midrule
\edit{Projection QP} & \edit{Up to five Clarabel attempts, each with at most
500 iterations. For primal objective $F$ and numerical dual bound $D$,
require $-10^{-7}(1+F)\le F-D\le2\times10^{-6}(1+F)$.}\\
\edit{Coherence} & \edit{Require $u\ge0$ and
$|p_i-(Su)_i|\le2\times10^{-6}+10^{-8}|(Su)_i|$.}\\
\edit{Nodewise protection} & \edit{Require $\Gamma_{R_i}(b_i,p_i)\le A_i$.
The stored binary values are converted exactly to rational numbers for
this test and allowance updates.}\\
\edit{Failure handling} & \edit{Unaccepted QP solutions use an SLSQP fallback;
projection exhaustion halts the trajectory. A failed release check triggers
contraction toward the reference.}\\\bottomrule
\end{tabular}\end{table}

\edit{The tolerances in Table~\ref{tab:numerics} are fixed numerical
settings. The objective gap checks how closely the solver approaches the projection
minimum. The separate rational-envelope test checks whether the forecast
satisfies each node's protection limit.}

\subsection{\edit{Release rules and state reset}}\label{sec:statepaths}
\edit{Cov and ML update their own allowances from their released predictions.
After repairing coordinates, Cov records the objective gap, while ML
also requires it to pass the acceptance threshold. On M5, Cov shrinks the
forecast and its representation together instead of clipping coordinates.
Section~\ref{sec:m5path} evaluates this numerical difference.}

Before the observation returns, Cov temporarily deducts the certified
amount $r_{t,i}=\Gamma_{R_i}(b_{t,i},p_{t,i})$ from the available allowance.
It restores this amount when it updates the allowance using the observed
error. With feedback before the next release,
\begin{equation}
(A_{t,i}-r_{t,i})+r_{t,i}+R_i(b_{t,i}-y_{t,i})^2-(p_{t,i}-y_{t,i})^2
=A_{t+1,i}.
\label{eq:reservation}
\end{equation}
The temporary deduction and restoration cancel, giving the same next
allowance as a direct update from the observed error. Cov and ML can
still have different allowances because they release different forecasts.

\begin{table}[htbp]\centering\small
\caption{State at the start of scoring.}\label{tab:reset}
\begin{tabular}{@{}p{.31\linewidth}p{.64\linewidth}@{}}\toprule
Quantity & Action\\\midrule
Protection allowance & Zero at the agreed start. Window tests restart for each window; the continuous path retains it for all 144 releases. The paired reset control clears only this allowance at release 73.\\
Expert histories & Retain returned raw errors, observations and ridge features.\\
Combination states & Retain Cov's recent expert errors and ML's sequential gain accumulators.\\
Variance weights and reference rule & Fixed before the corresponding calibration period. Continuous and reset-control paths share the early-window weights and each current reference.\\\bottomrule
\end{tabular}\end{table}

\FloatBarrier
\subsection{Backbone training and reuse}\label{sec:training}
Three independently initialized training runs use the same chronological
selection rules within each backbone. Table~\ref{tab:backbone_settings}
lists the fixed settings. Each PatchTST/DLinear calibration forecast is generated by a model
fitted and validated on earlier dates. DUET is fitted
before the earliest calibration period, except Labour's separately fitted
late stage. All postprocessors receive the same stored predictions from a given
backbone run. Continuous evaluation uses the early calibration period,
followed by both scoring windows. It does not repeat the late-window
calibration records, which already lie in the first scored half.

\begin{table}[htbp]\centering\footnotesize
\caption{Backbone settings.}\label{tab:backbone_settings}
\setlength{\tabcolsep}{3pt}
\begin{tabular}{llll}\toprule
Setting & PatchTST & DLinear & DUET\\
\midrule
Input / horizon & 24 / 12 & 24 / 12 & 24 / 12\\
Learning rate & $10^{-4}$ & $10^{-4}$ & $5\times10^{-4}$\\
Training loss & MAE & MAE & MAE + balancing loss\\
Training cap & 500 updates & 10,000 updates & 100 epochs\\
Early stopping & 50 checks & 50 checks & 5 epochs\\
Validation & 36 observations & 36 observations & 20\% prefix holdout\\
Normalization & RevIN & Identity & Author normalization\\
\bottomrule\end{tabular}
\end{table}
DUET adds the authors' expert-balancing loss to MAE with coefficient 1;
validation uses MAE alone. It uses one encoder layer, model width 64, feed-forward width 128,
two experts, one selected expert, one attention head and batch size 16.
Its configuration is fixed across datasets. Eleven original DUET windows are reused unchanged; the additional runs
cover the remaining initializations, Labour's early window and AEMO.
Separate fits on earlier dates supply forecast errors for training the
external adapters. Validation selects models using only pre-scoring
observations. Online Resource~2 (OR2) v11 lists the fitting dates and saved
training records.

The new AEMO short-window adapters use 216 historical forecast errors and
36 calibration observations at each window's boundary. The whole-month
adapters retain 108 support forecasts and 36 calibration observations.
On other datasets with shorter histories, adapter training uses only
the available earlier forecasts.

\FloatBarrier

\subsection{Observation definition and frozen boundaries}
The outcome is AEMO's estimated underlying demand: the sum of its
unadjusted half-hourly operational demand and measurement-based rooftop
PV generation estimate, in MW. Operational demand is read from
\texttt{OPERATIONAL\_DEMAND/ACTUAL}, field
\texttt{OPERATIONAL\_DEMAND}; PV is read from
\texttt{ROOFTOP/ACTUAL}, field \texttt{POWER}, with
\texttt{TYPE=MEASUREMENT}. AEMO defines this sum as estimated total
regional electricity consumption. The official sources are cited in
\MainSec{sec:datasets}. The resulting demand series estimates total regional consumption from
operational-demand measurements and estimated rooftop PV generation.

All five regions are retained in the order NSW, QLD, SA, TAS, VIC,
preceded by their system sum. The 2,928 half-hourly timestamps cover
1 September 2026 00:30 through 1 November 2026 00:00 in fixed AEST
(UTC+10). Both components contain every required region and timestamp. When
reports revise the same value, we use the latest \texttt{LASTCHANGED}
record; conflicting values tied at that time fail the completeness
check. We retain all observations without interpolation or deletion,
including both PV quality indicators, 0.7 and 1.

The original \texttt{TOTALDEMAND} registration is retained as an
ineligible signed-outcome attempt: 34 October half-hour values in SA
were negative. No forecasting method was scored on that attempt.
The replacement outcome was chosen from AEMO's physical definition.
Before acquiring these component data, the author saved the evaluation
protocol and 86 file hashes locally on 7 November 2026 at
03:26:17.918634 UTC. The hashes cover 84 Python implementation and
dependency files, one directory inventory and one protocol file.
The separate source manifest records the URLs and post-download hashes
of 520 original data ZIPs. These local timestamps, hashes and amendment records document when the
protocol was fixed. Related electricity fields had already been inspected;
the replacement outcome and its data acquisition followed this recorded
amendment.
The original registration and protocols remain in the historical v10 payload of Online Resource~2 v11.

The 520 raw ZIP files contain 3,594 CSV reports. The August trading-day archive supplies the records needed at the
September boundary, completing both component series. The operational-demand component retains
37 negative values (30 in October); adding the original PV estimate
yields nonnegative underlying demand at every region--time point.
October minima are 6,206, 5,445, 1,072, 910 and 4,313 MW in the five
regions, respectively. The system is their sum.

\begin{table}[htbp]\centering\footnotesize
\caption{AEMO boundaries and training runs.}\label{tab:aemoboundaries}
\setlength{\tabcolsep}{4pt}
\begin{tabular}{ll}\toprule
Stage & Setting\\\midrule
Backbone training / validation & Before 28 September 00:30\\
Support forecasts & 28 September 00:30--30 September 06:00 (108)\\
Calibration forecasts & 30 September 06:30--1 October 00:00 (36)\\
Scored month & 1 October 00:30--1 November 00:00 (1,488)\\
Training runs & Three per backbone; common validation-MAE rule\\
Input / base forecast horizon & 24 / 12 half-hours\\
Reference / allowance & NNLS / 0.5\%, $R=40401/40000$\\
Allowance reset & Zero only at the start of October\\
\bottomrule\end{tabular}
\end{table}
\subsection{Three backbones and three training runs}
PatchTST and DLinear whole-month predictions and release results are reused
unchanged. DUET adds three independently trained streams using the same
pre-support cutoff, 24/12 input/horizon and validation-MAE selection.
The shared monthly comparison contains 23 methods available on all three
backbones. Each protected method starts with zero once at October's first
release and settles all 1,488 outcomes without resetting.

All training runs are evaluated on the same finalized October
observations, so they measure sensitivity to initialization. The expanded
comparisons and short windows were added after the original new-outcome
study. Evaluation assumes feedback before the next release and uses
finalized data; it does not recreate historical reporting delays or
revisions.


\subsection{M5 publication-path sensitivity}\label{sec:m5path}
\edit{This sensitivity comparison changes only M5 Cov's numerical release
rule: joint contraction is replaced by the coordinate-repair rule used
on the other datasets. The two backbones, two windows and three runs
give twelve paired trajectories. Targets, reference, variance weights
and feedback order are fixed; each rule starts with zero allowance and
updates it from its own predictions.}

\edit{We compute RMSE from each trajectory's predictions and observations,
average the two windows within each backbone--run pair, and calculate
the relative reduction from joint contraction to coordinate repair.
Across the six resulting pairs, the mean reduction is
$4.62\times10^{-6}\%$ and the largest absolute change is
$2.93\times10^{-5}\%$. Both rules satisfy the cumulative limits.
The numerical release-rule difference has negligible aggregate accuracy
effect in this M5 comparison.}

\FloatBarrier
\Needspace{22\baselineskip}
\FloatBarrier

\section{Version map and numerical resources}\label{sec:resources}
OR2 v11 supplies the current unrounded per-run and grouped RMSEs,
paired percentages, dataset effects, training receipts and table checker.
Its numerical payload contains the stored contexts and releases for the
seven-dataset extension. The v10 archive preserves earlier analyses,
Labour provenance and original AEMO registrations unchanged.
The accompanying supplement-check addendum provides verified table
records and the corrected historical Li level summaries.

From the OR2 root, the aggregate command is
\begin{quote}\footnotesize\ttfamily
python current/check\_tables.py
\end{quote}
The optional stored-trajectory command extracts the supplied numerical
payload and checks its saved releases:
\begin{quote}\footnotesize\ttfamily
python current/check\_tables.py --trajectories
\end{quote}
The first command recomputes reported tables; the second also checks
the stored forecast sequences. Both use recorded predictions.
